\documentclass[11pt]{article}

\usepackage[margin=1in]{geometry}
\usepackage{graphicx}
\usepackage{booktabs}
\usepackage{amsmath,amssymb}
\usepackage{natbib}
\usepackage{hyperref}
\usepackage{microtype}
\usepackage{textgreek}

\bibliographystyle{unsrtnat}

\title{Design Principles for Inflammasome Inhibition by Pyrin-Only-Proteins}
\author{}
\date{}

\begin{document}

\maketitle

\begin{abstract}
Inflammasomes are filamentous cytosolic signaling platforms that are
indispensable for innate host defense, yet their persistent activity drives
autoinflammatory disease, cancer, and severe COVID-19. Endogenous
pyrin-only-proteins (POPs) are widely viewed as the principal throttles on
these platforms, with POP1, POP2, and POP3 assigned to ASC, NLRP2, and AIM2
respectively on the basis of 65--68\% sequence identity to each presumptive
target. How these small decoys actually recognize their substrates and stop
polymerization has never been tested directly, because prior work relied on
cellular assays that carry the full inflammasome cascade. Here we combine
Rosetta InterfaceAnalyzer energies on honeycomb sideviews of ASC, AIM2,
IFI16, NLRP3, and NLRP6 pyrin-domain (PYD) filaments with HEK293T
co-transfection imaging and in vitro FRET polymerization coupled to
negative-stain electron microscopy. POP1 is a poor inhibitor of ASC$_{PYD}$
($\leq 20$\% filament suppression at 1200~ng POP-eGFP and no detectable
kinetic effect up to 150~\textmu M) but potently suppresses AIM2, IFI16,
NLRP3, and NLRP6 PYD assembly (60--70\% in cells). POP2 and POP3 instead
block ASC$_{PYD}$ nucleation directly, extending the lag phase by up to
$\sim$15~min and abolishing mature filaments by electron microscopy. Interface
profiles reveal a simple design principle: inhibitory POPs must present both
a favourable recognition surface ($\Delta\Delta G \leq 3.5$) and at least one
strongly unfavourable repulsion interface ($\Delta\Delta G \geq 10.0$). These
results recast POPs as decoy adaptors and receptors whose targets are set by
interface complementarity rather than primary sequence, and supply a
protomer-level framework for therapeutic modulation of inflammasomes.
\end{abstract}

\section{Introduction}

Inflammasomes are multiprotein cytosolic signaling platforms that respond to
pathogen invasion, genotoxic stress, and a wide variety of danger signals,
and they represent a cornerstone of innate immune defense in mammals
\citep{broz2016inflammasomes, tartey2019inflammasomes}. A germline-encoded
pattern-recognition receptor, such as the AIM2-like receptors (ALRs) AIM2
and IFI16 or the NOD-like receptors (NLRs) NLRP3 and NLRP6, nucleates
filamentous polymers through its pyrin domain (PYD), which in turn seed
PYD filaments of the bipartite adaptor ASC. The ASC filament then templates
caspase-1 activation via CARD--CARD filaments, driving IL-1$\beta$ and IL-18
maturation and pyroptotic cell death
\citep{lu2014unified, hornung2009aim2, fernandesalnemri2009aim2,
roberts2009hin200}. This sequential, prion-like assembly is exceptionally
robust: each step amplifies the incoming signal, the polymerized state
persists once nucleated, and the response is effectively binary and
irreversible, properties ideally suited for rapid host defense
\citep{cai2014prion, franklin2014adaptor, matyszewski2018digital}. The very
same properties, however, render persistent inflammasome activity a liability.
Dysregulated assembly has been implicated in cryopyrin-associated periodic
syndromes and other autoinflammatory conditions, in cancer-associated
inflammation, and in severe COVID-19 driven by monocyte inflammasome
activation \citep{sharma2018pyrin, vora2021inflammasome, junqueira2022sarscov2}.
Extrinsic regulators that can throttle inflammasomes without shutting down
innate immunity are therefore biologically essential.

A major class of such endogenous regulators are the mammalian
pyrin-only-proteins (POPs), small six-helix bundles of roughly 100 residues
that bear a PYD but lack the downstream effector domains present in sensor
or adaptor proteins \citep{indramohan2018cops}. Three human POPs have been
characterized to date. POP1 is 65\% identical to the PYD of ASC and has been
proposed to block ASC nucleation directly \citep{dealmeida2015pyrin}. POP2
is 68\% identical to the PYD of NLRP2 and has been implicated in inhibiting
ASC-dependent inflammasomes and AIM2 \citep{periasamy2017pop2,
ratsimandresy2017pop2}. POP3 is 67\% identical to the PYD of AIM2 and has
been reported to restrict AIM2- and IFI16-driven responses to viral dsDNA
\citep{khare2014pop3}. This tripartite, homology-driven assignment forms
the current textbook model of POP target specificity.

These assignments rest, however, on indirect inference. Cellular and in vivo
measurements of POP activity necessarily confound direct PYD--PYD recognition
with the many parallel interactions of a full inflammasome cascade, and no
biochemical assay has yet resolved whether a given POP blocks nucleation,
elongation, or both against a defined PYD substrate. The premise that primary
sequence identity predicts POP targets is moreover difficult to reconcile
with the architecture of PYD filaments. Structural work on ASC, AIM2,
NLRP3, and NLRP6 filaments has shown that every PYD engages its neighbours
through six largely distinct protomer interfaces of Type 1, Type 2, and
Type 3 character, and that receptor PYDs with essentially no pairwise
sequence homology all converge on the single adaptor ASC$_{PYD}$
\citep{lu2014unified, shen2019nlrp6, hochheiser2022directionality,
matyszewski2021distinct, morrone2015assembly, sharif2019nek7,
andreeva2021nlrp3}. This degenerate, code-like recognition implies that
local interface complementarity, rather than global sequence identity,
sets specificity.

To test the sequence-homology model directly and to ask how favourable and
unfavourable protomer-level interactions combine to permit or prevent PYD
filament assembly, we combined three orthogonal readouts. First, Rosetta
InterfaceAnalyzer computations on honeycomb-like sideviews of ASC, AIM2,
IFI16, NLRP3, and NLRP6 PYD filaments quantify the interface energies
($\Delta G$, in Rosetta Energy Units) for homotypic PYD--PYD contacts and
for all hypothetical POP--PYD substitutions, using cryo-EM and crystal
structures of each filament or its closest structural homologue
\citep{lu2014unified, shen2019nlrp6, hochheiser2022directionality,
sharif2019nek7, matyszewski2021distinct}. Second, HEK293T cells, which lack
endogenous inflammasome components and POPs, provide a clean co-transfection
context in which mCherry-tagged PYDs and eGFP-tagged POPs can be scored for
filament or punctum suppression in isolation. Third, FRET-based polymerization
kinetics using purified recombinant POPs and PYDs, complemented by
negative-stain electron microscopy (nsEM), dissect whether a given POP blocks
the rate-limiting nucleation step, attenuates elongation, or both
\citep{mazanek2019fret, matyszewski2018digital, morrone2015assembly}.

The central findings of this work are threefold. First, POP1 is a poor
inhibitor of ASC$_{PYD}$, suppressing filament assembly by no more than
$\sim$20\% in HEK293T cells at 1200~ng POP-eGFP and showing no detectable
effect on in vitro polymerization up to 150~\textmu M; this overturns the
canonical POP1-as-ASC-inhibitor model. Second, POP1 instead acts
predominantly on upstream receptor PYDs (AIM2, IFI16, NLRP3, NLRP6),
suppressing their filament assembly by 60--70\% in cells and interfering
with elongation in vitro. Third, POP2 and POP3 potently block ASC$_{PYD}$
nucleation, extending the lag phase by up to $\sim$15~min in FRET
experiments and abolishing or strongly reducing ASC$_{PYD}$ filaments in
nsEM. Synthesising these observations with the Rosetta interface analyses,
we put forward a design principle in which inhibitory POPs must
simultaneously present at least one favourable recognition interface and
at least one strongly unfavourable repulsion interface with their PYD
substrate. Together these results redefine POPs as decoy adaptors and
receptors, demote sequence homology as a predictor of function, and provide
a protomer-interface framework that can guide the rational design of POP
variants or small-molecule mimics for therapeutic modulation of persistent
inflammasome activity.

\section{Materials and Methods}

\subsection{Rosetta interface analysis}

Interface energies ($\Delta G$, in Rosetta Energy Units) were computed
using Rosetta \texttt{InterfaceAnalyzer} on honeycomb-like sideviews of PYD
filaments. In each honeycomb, the central protomer engages six neighbours
through Type 1a, Type 1b, Type 2a, Type 2b, Type 3a, and Type 3b interfaces
(Figure~\ref{fig:honeycomb}), and the total interface energy is the sum of
all six. Homotypic filaments were built from the deposited cryo-EM and
crystal structures of ASC$_{PYD}$ (PDB 3J63) \citep{lu2014unified},
AIM2$_{PYD}$ (PDB 7K3R) \citep{matyszewski2021distinct}, NLRP3$_{PYD}$
(PDB 7PDZ) \citep{hochheiser2022directionality}, and NLRP6$_{PYD}$
(PDB 6NCV) \citep{shen2019nlrp6}. IFI16$_{PYD}$ was modeled onto the
eGFP-AIM2 filament (PDB 6MB2) given the absence of a direct structure.
POP1 (PDB 4QOB) was taken directly; POP2 and POP3 were built by homology
modeling onto NLRP3$_{PYD}$ (PDB 7PDZ) and AIM2$_{PYD}$ (PDB 7K3R)
respectively. For each POP--PYD pair, the central PYD in the honeycomb was
substituted with the POP and the six protomer interfaces were re-evaluated,
keeping the rest of the filament fixed. Interactions with
$\Delta\Delta G \leq 3.5$ (POP--PYD minus homotypic PYD--PYD) were scored
as favourable recognition surfaces, and those with $\Delta\Delta G \geq 10.0$
as strongly unfavourable repulsion interfaces.

\subsection{Cell culture and co-transfection imaging}

HEK293T cells (ATCC, CRL-11268) were seeded into 12-well plates
($0.1 \times 10^{6}$ cells per well) containing round cover glass
(20~mm) at 70\% confluence. POP-eGFP or eGFP alone plasmids (600 or
1200~ng) were co-transfected with PYD-mCherry plasmids (300 or 600~ng)
using Lipofectamine 2000 (Invivogen). After 16~h, cells were washed twice
with 1$\times$ phosphate-buffered saline (PBS), fixed with 4\%
paraformaldehyde, and mounted on glass slides. DAPI was used to stain
nuclei. Images were acquired on a Cytation 5 imager and analyzed with
Gen5 software. All cellular suppression values reported here are
computed from $n \geq 4$ independent fields relative to the eGFP-only
control.

\subsection{Recombinant protein expression and purification}

POP1 was cloned into pET21b and POP2 and POP3 into pET28b with an
N-terminal His$_6$-MBP tag followed by a TEVp cleavage site. All
recombinant proteins were expressed in \emph{Escherichia coli} BL21
Rosetta2(DE3) and purified by Ni$^{2+}$-NTA followed by size-exclusion
chromatography (Superdex 75, Cytiva) in storage buffer
(40~mM HEPES-NaOH pH 7.4, 400~mM NaCl, 2~mM dithiothreitol,
0.5~mM EDTA, 10\% glycerol). Proteins were stored at $-80$~$^{\circ}$C.
Monomeric POP1 and MBP-POP2/3 peaks were collected and used in downstream
assays; untagged POP2/3 were generated by TEVp cleavage of the MBP tag
(5~\textmu M TEVp, 30~min).

\subsection{FRET polymerization kinetics}

FRET assays used donor- and acceptor-labeled MBP-PYD mixed at a 1:1 ratio
to a final PYD concentration of 2.5~\textmu M. Polymerization was triggered
by the addition of TEVp, which cleaves the MBP tag and releases the PYD for
self-assembly. Increasing concentrations of each POP were included at the
start of the reaction. For ASC$_{PYD}$, POP1 was tested at 50 and 150~\textmu M,
POP2 at 3.3, 6.7, 13.3, 26.7 and 40~\textmu M, and POP3 at 3.25, 7.5, 15, 20
and 30~\textmu M. For AIM2$_{PYD}$, POP1 was tested at 25, 50, 100 and
150~\textmu M, POP2 at 12.5, 25, 50 and 75~\textmu M, and POP3 at 7.5, 15,
30, 40 and 50~\textmu M. For NLRP6$_{PYD}$, POP1 was tested at 25, 50 and
100~\textmu M, POP2 at 30, 60 and 120~\textmu M, and POP3 at 15, 30 and
60~\textmu M. Half-times for polymerization ($t_{1/2}$) and IC$_{50}$ values
(defined as the POP concentration that decreases $1/t_{1/2}$ by 50\%) were
computed as in prior work \citep{matyszewski2018digital,
matyszewski2021distinct}. All values reported are averages of at least
three independent measurements ($n = 3$).

\subsection{Full-length dsDNA binding assays}

Fluorescence anisotropy (FA)-based binding of a 60-bp FAM-labeled dsDNA
(10~nM) to AIM2$_{FL}$ (100~nM) or IFI16$_{FL}$ (200~nM) was measured with
increasing concentrations of POPs. The MBP tag was pre-cleaved by TEVp for
30~min before addition to the binding reaction. IC$_{50}$ values were
computed as previously described \citep{matyszewski2018digital,
matyszewski2021distinct}.

\subsection{Negative-stain electron microscopy}

For nsEM, TEVp-cleaved MBP-PYDs were allowed to polymerize in the
presence or absence of POPs. In ASC$_{PYD}$ experiments POP1 was used at
150~\textmu M, POP2 at 40~\textmu M, and POP3 at 30~\textmu M. Samples
were adsorbed onto carbon-coated copper grids and imaged as previously
described \citep{shen2019nlrp6, matyszewski2021distinct}.

\section{Results}

\subsection{Rosetta interface analyses suggest broad POP target
specificities that exceed sequence-homology predictions}

Because the existing sequence-homology model of POP specificity assumes
that primary identity predicts filament-level recognition, we first asked
whether computed interface energies bore this out. We carried out Rosetta
InterfaceAnalyzer calculations on honeycomb sideviews of ASC, AIM2, IFI16,
NLRP3, and NLRP6 PYD filaments (Figure~\ref{fig:honeycomb}), systematically
substituting the central protomer with POP1, POP2, or POP3, and
quantifying interface energies for each of the six protomer contacts and
for their top-half and bottom-half sums. Homotypic PYD--PYD interactions
were highly symmetric, with Type 1a and Type 1b interfaces essentially
degenerate (for example, ASC$_{PYD}$--ASC$_{PYD}$ at $-22.1$~REU), which
is consistent with bidirectional filament growth and with the degenerate
assembly geometry reported by prior cryo-EM studies
\citep{lu2014unified, matyszewski2021distinct, hochheiser2022directionality}.

A clear quantitative hierarchy emerged. Homotypic $\Delta G$ values were
always more favourable than any POP--PYD interaction on the corresponding
filament: for ASC, the homotypic top-half sum reached $-35.9$~REU while
POP1--ASC$_{PYD}$ achieved only $-24.7$~REU (Figure~\ref{fig:rosetta_dg}).
This energetic gap immediately implies that thermodynamics alone will
favour self-assembly over POP engagement, and that inhibition requires
super-stoichiometric POP concentrations. Even so, the POP--PYD energies
spread surprisingly broadly. POP1 showed the most favourable top-half
interface with ASC$_{PYD}$ but POP3 matched it at the bottom half
(POP1--ASC$_{PYD}$ top $=-34.3$~REU vs POP3--ASC$_{PYD}$ bottom
$=-34.6$~REU), indicating that a POP can in principle dock onto the ASC
filament through more than one protomer face. For AIM2$_{PYD}$ all three
POPs were comparable on the bottom half; POP3 was the most favourable POP
for IFI16$_{PYD}$ with a top-half $\Delta G$ of $-9.6$~REU, yet even POP1
and POP2 exceeded homotypic IFI16--IFI16 energies at individual
interfaces. Most striking, for NLRP3$_{PYD}$, POP1 was nearly as
favourable as homotypic contacts (POP1--NLRP3$_{PYD}$ $=-30.5$~REU
vs NLRP3$_{PYD}$--NLRP3$_{PYD}$ $=-31.4$~REU), while POP2 and POP3
interactions with NLRP3$_{PYD}$ were much less favourable
(Figure~\ref{fig:rosetta_dg}).

Taken together, these in silico analyses are inconsistent with the
sequence-homology prediction that each POP has a single cognate PYD
target. They rationalize portions of the prior model -- POP1 does indeed
score best against ASC$_{PYD}$ on a per-interface basis -- but they also
imply that POP1 should be even more effective against NLRP3$_{PYD}$ and
that POP2 and POP3 should engage AIM2$_{PYD}$ and IFI16$_{PYD}$ at least
as well as POP3 does. Since this is difficult to reconcile with in vivo
phenotypes attributing primary ASC inhibition to POP1, we turned to
direct biochemical and cellular experiments to resolve the discrepancy.

\begin{figure}[htbp]
\centering
\includegraphics[width=0.55\textwidth]{figures/fig_honeycomb_schematic.pdf}
\caption{Honeycomb sideview strategy used for Rosetta interface analysis.
The central PYD protomer engages six neighbours through Type 1a/1b, Type
2a/2b and Type 3a/3b interfaces, which together define the interfaces
probed in all homotypic and POP--PYD calculations.}
\label{fig:honeycomb}
\end{figure}

\begin{figure}[htbp]
\centering
\includegraphics[width=0.78\textwidth]{figures/fig_rosetta_dg.pdf}
\caption{Rosetta InterfaceAnalyzer energies ($\Delta G$, Rosetta Energy
Units) for the top-half interfaces of homotypic PYD--PYD filaments and
POP--PYD substitutions. ASC$_{PYD}$--ASC$_{PYD}$ reaches $-35.9$~REU while
POP1--ASC$_{PYD}$ is only $-24.7$~REU; POP1--NLRP3$_{PYD}$ ($-30.5$~REU)
is nearly isoenergetic with NLRP3$_{PYD}$--NLRP3$_{PYD}$ ($-31.4$~REU);
POP3--IFI16$_{PYD}$ (top half, $-9.6$~REU) is the most favourable POP
interaction with IFI16$_{PYD}$.}
\label{fig:rosetta_dg}
\end{figure}

\subsection{Mechanisms of ASC inhibition by POPs: POP1 is a poor ASC
inhibitor, while POP2 and POP3 directly block ASC nucleation}

To test POP specificity in a clean cellular context, we used HEK293T cells,
which lack endogenous inflammasome components and POPs
\citep{hornung2009aim2, fernandesalnemri2009aim2}. ASC$_{PYD}$-mCherry
(300~ng) or ASC$_{FL}$-mCherry (300~ng) was co-transfected with POP-eGFP
or eGFP alone at 600 or 1200~ng, and filaments or puncta were scored after
16~h of expression. Surprisingly, and in direct contradiction to the
classical POP1-as-ASC-inhibitor model, POP1-eGFP minimally inhibited
ASC$_{PYD}$-mCherry filament formation, with $\leq 20$\% suppression at
1200~ng relative to the eGFP control. POP2 and POP3, by contrast, strongly
reduced ASC$_{PYD}$ filaments: mCherry signal became diffuse and the
remaining structures were short and non-linear. A parallel effect was
seen on ASC$_{FL}$ puncta, where POP1 reduced their number by only
$\sim$30\% while POP2 and POP3 achieved up to $\sim$80\% reduction.

To dissect the step at which POP2 and POP3 act, we developed a recombinant
FRET-based polymerization assay. Recombinant POP1 behaved as a monomer in
size-exclusion chromatography, whereas POP2 and POP3 required the
N-terminal MBP tag for solubility and, upon TEVp cleavage, formed
elongated oligomers visualized by nsEM. Polymerization of donor- and
acceptor-labeled ASC$_{PYD}$ (2.5~\textmu M, donor:acceptor 1:1) triggered
by TEVp cleavage of MBP proceeded through the two expected phases: a
rate-limiting nucleation lag followed by rapid elongation
\citep{matyszewski2018digital, mazanek2019fret}. POP1 up to 150~\textmu M
had no detectable effect on either lag or elongation phase and did not
alter nsEM filament formation, consistent with its negligible activity in
cells. POP2 (at 3.3, 6.7, 13.3, 26.7, and 40~\textmu M) and POP3 (at 3.25,
7.5, 15, 20, and 30~\textmu M), however, prolonged the ASC$_{PYD}$ lag
phase in a dose-dependent manner by up to $\sim$15~min and additionally
reduced the elongation slope by $\leq 20$\%
(Figure~\ref{fig:fret_concs}).

The nsEM images were more stringent still. In the presence of POP2, no
ASC$_{PYD}$ filaments were observed; with POP3, only a few short
filaments formed. Because the FRET signal nonetheless rose during the
lag-extended regime, POPs must trap transient pseudo-nucleation
intermediates that do not mature into stable filaments. Together these
results reassign POP1 out of the ASC-inhibitor category and identify POP2
and POP3 as dose-dependent nucleation blockers that physically prevent
ASC$_{PYD}$ filament maturation.

\begin{figure}[htbp]
\centering
\includegraphics[width=\textwidth]{figures/fig_fret_concentrations.pdf}
\caption{POP concentration series used in FRET-based polymerization
assays against ASC$_{PYD}$, AIM2$_{PYD}$ and NLRP6$_{PYD}$ (each at
2.5~\textmu M). POP1 is tested up to 150~\textmu M against ASC$_{PYD}$
and AIM2$_{PYD}$ and up to 100~\textmu M against NLRP6$_{PYD}$; POP2 and
POP3 operate at lower absolute concentrations.}
\label{fig:fret_concs}
\end{figure}

\subsection{Re-examining the Rosetta analyses in light of biochemical
experiments reveals a favourable-plus-unfavourable interface rule}

Having found that POP1 fails to inhibit ASC$_{PYD}$ despite scoring
favourably in bulk Rosetta energies, we asked whether a closer
inspection of individual protomer interfaces would rationalize the
discrepancy. To do so, we decomposed the six interfaces per honeycomb
and computed the per-interface $\Delta\Delta G$ (POP--PYD minus
homotypic PYD--PYD) for each POP--ASC pair, reasoning that filament
assembly is a six-fold interaction and that bulk $\Delta G$ sums may
mask individual outliers that matter for polymerization.

The per-interface decomposition was striking in its simplicity.
POP2--ASC and POP3--ASC profiles each contained multiple strongly
unfavourable interfaces, with $\Delta\Delta G \geq 10.0$~REU, alongside
several favourable recognition surfaces ($\Delta\Delta G \leq 3.5$~REU).
POP1--ASC, by contrast, contained favourable recognition surfaces but
lacked any strongly unfavourable repulsion interface. In other words,
POP1 can dock onto ASC$_{PYD}$ but does not present a surface that is
incompatible with the next protomer in the filament, so assembly can
still propagate around the docked POP. POP2 and POP3, by contrast,
combine recognition and repulsion on the same molecule and therefore
interrupt the filament once they are bound.

This interface decomposition crystallised a concrete hypothesis:
inhibitory POPs must simultaneously present a favourable recognition
surface, which permits them to dock onto a target PYD, and at least one
strongly unfavourable interface, which prevents a subsequent protomer
from fully engaging the POP and propagates a defect along the filament.
Because filament assembly requires all six interfaces to be mutually
compatible, even a single unfavourable surface is sufficient to block
polymerization once the POP is bound. We term this the
favourable-plus-unfavourable interface rule and we test it below
against the ALR and NLR substrates, where it provides quantitative
predictions that were not available from sequence homology alone.

\subsection{Mechanisms of ALR inhibition by POPs}

Inflammasomes of the ALR family respond to cytosolic dsDNA
\citep{hornung2009aim2, fernandesalnemri2009aim2, roberts2009hin200,
morrone2015assembly} and are directly clinically relevant as
autoantigens in systemic lupus erythematosus
\citep{antiochos2022dna, kerur2011ifi16}, so we next asked whether the
favourable-plus-unfavourable interface rule would also account for POP
action against AIM2 and IFI16. The Rosetta calculations were
encouraging on this front: AIM2$_{PYD}$ and IFI16$_{PYD}$ interface
profiles contained a mixture of favourable and unfavourable interfaces
with all three POPs, matching the pattern we had associated with
inhibitory activity against ASC.

Cellular imaging in HEK293T recapitulated this prediction almost
perfectly. Compared with co-transfection of eGFP alone, POP1-eGFP
reduced AIM2$_{PYD}$ and IFI16$_{PYD}$ filament counts far more
effectively than it reduced ASC$_{PYD}$ filaments: at 1200~ng POP1,
AIM2$_{PYD}$ and IFI16$_{PYD}$ assemblies were inhibited $\sim$60\%,
compared with only $\sim$20\% for ASC$_{PYD}$. POP2 and POP3 essentially
eliminated AIM2$_{PYD}$ and IFI16$_{PYD}$ filaments in cells. The same
hierarchy held when the POP dose was reduced: at 600~ng POP3,
AIM2$_{PYD}$ assembly was suppressed by 70\% while ASC$_{PYD}$ and
NLR PYDs were reduced by only 40--50\%, underscoring the intrinsic
preference of POP3 for ALR substrates even at concentrations where its
action on ASC is muted.

In vitro FRET polymerization kinetics then identified the mechanistic
step at which this suppression operates. AIM2$_{PYD}$ (2.5~\textmu M)
was assayed with POP1 (25, 50, 100, and 150~\textmu M), POP2 (12.5, 25,
50, and 75~\textmu M), and POP3 (7.5, 15, 30, 40, and 50~\textmu M).
In striking contrast to their action on ASC$_{PYD}$, all three POPs
reduced the elongation slope of AIM2$_{PYD}$ in a dose-dependent manner
with only minimal change to the lag phase, and POP3 displayed the lowest
IC$_{50}$ (Figure~\ref{fig:fret_concs}). nsEM imaging confirmed the
kinetic reading: POP1 made AIM2$_{PYD}$ filaments both shorter and
less numerous, whereas POP2 and POP3 abolished filament formation
entirely. When the full-length proteins were assayed by fluorescence
anisotropy, AIM2$_{FL}$ dsDNA binding was unaffected by POP1 but
inhibited by POP2 and POP3, while for IFI16$_{FL}$ only POP3 was
inhibitory. All three POPs can therefore interfere with ALR filament
assembly, with POP3 the strongest inhibitor on a per-concentration
basis, and the presence of the activating ligand dsDNA partially
rescues assembly by stabilising homotypic receptor--receptor
recognition. Crucially, POPs attack ALRs at the elongation step rather
than at nucleation, which is the opposite regime to their action on
ASC and establishes that the same POP can operate in distinct kinetic
modes depending on which favourable and unfavourable interfaces it
presents to the substrate.

\subsection{POP1 acts as a decoy adaptor for NLR receptor PYDs}

Because POP1 fails as an ASC inhibitor yet engages AIM2, IFI16, NLRP3,
and NLRP6 PYDs in Rosetta calculations and cellular assays, we
hypothesized that POP1 is better understood as a decoy adaptor that
intercepts upstream receptors before they can engage ASC
\citep{lu2014unified}. Rosetta analyses supported this view. POP1 showed
a mixture of favourable and strongly unfavourable interfaces with
NLRP3$_{PYD}$ and NLRP6$_{PYD}$, exactly the pattern predicted by our
favourable-plus-unfavourable rule. POP2 and POP3 showed largely
unfavourable $\Delta G$ values on these substrates but retained at least
one recognition-capable interface (Type 3a, $\Delta G \sim -9$) that
could serve as an anchor.

Cellular imaging in HEK293T delivered a sharp quantitative picture of
POP1's revised role. NLRP2$_{PYD}$-mCherry failed to form filaments in
HEK293T cells, precluding a direct test of POP2 against its presumed
cognate target despite high sequence similarity; this failure of
NLRP2$_{PYD}$ to polymerize in isolation is itself informative, and it
further emphasizes that sequence identity alone is not a sufficient
design criterion. Against NLRP3$_{PYD}$ and NLRP6$_{PYD}$, POP1 was
highly effective: at 1200~ng, POP1 suppressed NLRP3$_{PYD}$ assembly by
70\% and NLRP6$_{PYD}$ assembly by 60\%, compared with only $\sim$20\%
for ASC$_{PYD}$ (Figure~\ref{fig:pop1_suppression}). The dose and target
hierarchy were also informative. At 600~ng POP2, ASC$_{PYD}$ assembly
was abolished, but NLRP3$_{PYD}$ and NLRP6$_{PYD}$ were barely
suppressed, and at 600~ng POP3, AIM2$_{PYD}$ was suppressed by 70\%
while ASC$_{PYD}$, NLRP3$_{PYD}$ and NLRP6$_{PYD}$ were reduced by only
40--50\%.

\begin{figure}[htbp]
\centering
\includegraphics[width=0.78\textwidth]{figures/fig_pop1_suppression.pdf}
\caption{Cellular suppression of PYD filament assembly by 1200~ng
POP1-eGFP in HEK293T cells co-transfected with 300~ng of the indicated
mCherry-tagged PYD. POP1 is a weak inhibitor of ASC$_{PYD}$ ($\sim$20\%
suppression) but suppresses upstream receptor PYDs by 60--70\%, consistent
with POP1 acting as a decoy adaptor upstream of ASC.}
\label{fig:pop1_suppression}
\end{figure}

In vitro FRET kinetics further dissected POP action against
NLRP6$_{PYD}$ (2.5~\textmu M). POP1 at 25, 50, and 100~\textmu M
predominantly inhibited elongation; POP2 at 30, 60, and 120~\textmu M
inhibited nucleation; and POP3 at 15, 30, and 60~\textmu M reduced both.
nsEM images confirmed that POP1, POP2, and POP3 each reduced the number
and length of NLRP6$_{PYD}$ filaments. Consistent with the pseudo-nucleus
trap observed against ASC$_{PYD}$, oligomers did form (FRET signal
increased) but did not mature into intact filaments. POP1 thus behaves
functionally as a decoy ASC that interferes with upstream PYDs rather
than with ASC itself, and a simultaneous combination of favourable and
unfavourable interfaces is both necessary and sufficient to block PYD
filament assembly (Figure~\ref{fig:model}).

\begin{figure}[htbp]
\centering
\includegraphics[width=0.95\textwidth]{figures/fig_model_schematic.pdf}
\caption{Canonical versus revised model of POP1 function.
Sequence-homology based models placed POP1 as a direct ASC inhibitor
(top). Our combined Rosetta, cellular, and biochemical data instead
indicate that POP1 acts as a decoy adaptor that throttles upstream
receptor PYDs, while POP2 and POP3 additionally block ASC nucleation
(bottom).}
\label{fig:model}
\end{figure}

\subsection{Summary tables}

Table~\ref{tab:rosetta} consolidates the principal Rosetta interface
energies used in this analysis, Table~\ref{tab:cell} the cellular
suppression values at 1200~ng POP-eGFP, and Table~\ref{tab:kinetic} the
kinetic regimes identified by FRET.

\begin{table}[htbp]
\centering
\caption{Selected Rosetta interface energies ($\Delta G$, REU) for
homotypic PYD--PYD filaments and POP--PYD interactions. Homotypic
$\Delta G$ values are more favourable than any POP--PYD interaction on
the same filament.}
\label{tab:rosetta}
\begin{tabular}{llrr}
\toprule
PYD target & Partner & Top-half $\Delta G$ (REU) & Bottom-half
$\Delta G$ (REU) \\
\midrule
ASC$_{PYD}$   & ASC$_{PYD}$ (homotypic)   & $-35.9$ & symmetric \\
ASC$_{PYD}$   & POP1                      & $-24.7$ & $-34.3$ \\
ASC$_{PYD}$   & POP3                      & --      & $-34.6$ \\
NLRP3$_{PYD}$ & NLRP3$_{PYD}$ (homotypic) & $-31.4$ & symmetric \\
NLRP3$_{PYD}$ & POP1                      & $-30.5$ & -- \\
IFI16$_{PYD}$ & POP3                      & $-9.6$  & -- \\
\bottomrule
\end{tabular}
\vspace{0.4em}
\\ Exemplar symmetric Type 1a/1b ASC$_{PYD}$--ASC$_{PYD}$ interface: $\Delta G = -22.1$~REU.
\end{table}

\begin{table}[htbp]
\centering
\caption{Cellular suppression of PYD filament and puncta formation by
POP-eGFP (1200~ng) relative to eGFP alone in HEK293T cells co-transfected
with 300~ng PYD-mCherry ($n \geq 4$). Values are reported as approximate
percent reductions taken verbatim from the experimental log.}
\label{tab:cell}
\begin{tabular}{lccc}
\toprule
PYD target (mCherry) & POP1 (\%) & POP2 (\%) & POP3 (\%) \\
\midrule
ASC$_{PYD}$    & $\leq 20$ & up to $\sim 80$ & up to $\sim 80$ \\
ASC$_{FL}$ puncta & $\sim 30$ & up to $\sim 80$ & up to $\sim 80$ \\
AIM2$_{PYD}$   & $\sim 60$ & essentially eliminated & essentially eliminated \\
IFI16$_{PYD}$  & $\sim 60$ & essentially eliminated & essentially eliminated \\
NLRP3$_{PYD}$  & $70$      & barely suppressed at 600~ng & 40--50 at 600~ng \\
NLRP6$_{PYD}$  & $60$      & barely suppressed at 600~ng & 40--50 at 600~ng \\
\bottomrule
\end{tabular}
\end{table}

\begin{table}[htbp]
\centering
\caption{Kinetic regime of PYD polymerization inhibition by POPs in
FRET assays (2.5~\textmu M PYD). ``Nucleation'' indicates an extension
of the lag phase; ``Elongation'' indicates a reduction in the exponential
or linear slope.}
\label{tab:kinetic}
\begin{tabular}{llllc}
\toprule
PYD target & POP & Nucleation effect & Elongation effect & POP range (\textmu M) \\
\midrule
ASC$_{PYD}$   & POP1 & none up to 150~\textmu M & none & 50, 150 \\
ASC$_{PYD}$   & POP2 & extended up to $\sim 15$~min & $\leq 20$\% slope reduction & 3.3--40 \\
ASC$_{PYD}$   & POP3 & extended up to $\sim 15$~min & $\leq 20$\% slope reduction & 3.25--30 \\
AIM2$_{PYD}$  & POP1 & minimal & dose-dependent reduction & 25--150 \\
AIM2$_{PYD}$  & POP2 & minimal & dose-dependent reduction & 12.5--75 \\
AIM2$_{PYD}$  & POP3 & minimal & dose-dependent reduction (lowest IC$_{50}$) & 7.5--50 \\
NLRP6$_{PYD}$ & POP1 & minor & predominant & 25--100 \\
NLRP6$_{PYD}$ & POP2 & predominant & minor & 30--120 \\
NLRP6$_{PYD}$ & POP3 & reduced & reduced & 15--60 \\
\bottomrule
\end{tabular}
\end{table}

\section{Discussion}

\subsection{Redefining the target specificity and inhibitory mechanisms of
POPs}

Inflammasomes are exceptionally stable, prion-like polymers whose aberrant
activity has been tied to autoimmune disease, cancer-associated
inflammation, and severe COVID-19
\citep{cai2014prion, franklin2014adaptor, sharma2018pyrin,
vora2021inflammasome, junqueira2022sarscov2}. Regulation must therefore
act either before signal amplification occurs, by suppressing nucleation,
or early in elongation, by preventing additional protomers from joining.
Our data show that POPs exploit both regimes but do so without
co-assembling into the mature filament, unlike CARD-only proteins that
are known to co-assemble with the caspase-1 CARD
\citep{indramohan2018cops}. Excess POP stoichiometry is required, because
homotypic PYD--PYD interactions always score as the most favourable in
Rosetta calculations -- a point that is especially sharp when an
activating ligand such as dsDNA further stabilises receptor assembly
\citep{morrone2015assembly}. POPs therefore appear designed to modulate,
rather than terminate, innate immune signaling, which is consistent with
their dispensability for basal host defense and their induction by
cytokine signaling \citep{dealmeida2015pyrin, ratsimandresy2017pop2,
khare2014pop3}.

\subsection{Revised role of POP1 as a decoy ASC}

Our most unexpected finding is that POP1 is a poor inhibitor of
ASC$_{PYD}$. In HEK293T cells, 1200~ng POP1-eGFP suppressed ASC$_{PYD}$
filaments by no more than $\sim$20\%, and in vitro POP1 up to
150~\textmu M left ASC$_{PYD}$ polymerization kinetics and nsEM
filament formation unchanged. Yet POP1 was clearly effective against
AIM2$_{PYD}$ and IFI16$_{PYD}$ ($\sim$60\% suppression in cells) and even
more so against NLRP3$_{PYD}$ (70\%) and NLRP6$_{PYD}$ (60\%). Because
POP1 expression is induced by IL-1$\beta$ and by TLR/IL-1R signaling
\citep{dealmeida2015pyrin}, it is best understood not as an ASC brake
but as a negative-feedback regulator that sequesters upstream receptors.
This recasts POP1 as a functional ASC mimic -- a decoy adaptor -- whose
physiological role is to stabilise rather than terminate inflammasome
activation by biasing upstream receptors away from productive ASC
engagement. Figure~\ref{fig:model} summarises this revised assignment.

\subsection{POP2 and POP3 as dual-target inhibitors}

POP2 and POP3 behave quite differently. Both potently suppress
ASC$_{PYD}$ nucleation -- extending the lag phase by up to $\sim$15~min
in FRET and essentially abolishing ASC$_{PYD}$ filaments in nsEM -- and
both also interfere with AIM2$_{PYD}$, IFI16$_{PYD}$, and NLR filament
assemblies. POP2 is induced by both pro- and anti-inflammatory cytokines
and here functions as a pan-inflammasome inhibitor
\citep{periasamy2017pop2, ratsimandresy2017pop2}, while POP3 is induced
exclusively by interferons and therefore operates in a strongly
ALR-biased physiological context \citep{khare2014pop3}. Our data make
clear that both POP2 and POP3 are dual-target, dual-regime inhibitors:
they block nucleation where the target is ASC and elongation where the
target is AIM2, and they combine these effects for NLR substrates.
Sequence homology is not sufficient to predict either the target or the
kinetic regime; POP2 is only moderately effective against its
sequence-cognate NLRP2$_{PYD}$ (which in our hands failed to polymerize
in HEK293T) and strongly active against the sequence-distant
ASC$_{PYD}$.

\subsection{Design principles and the favourable-plus-unfavourable
interface rule}

The interface-level decomposition of the Rosetta calculations yields a
crisp structural rule for PYD filament inhibition. POP1--ASC$_{PYD}$
lacks any interface with $\Delta\Delta G \geq 10.0$ and fails to inhibit
polymerization despite possessing favourable recognition surfaces.
POP1--NLRP3$_{PYD}$ has only one strongly unfavourable interface yet is
inhibitory. POP2/3--NLRP6$_{PYD}$ lack strongly favourable interactions
overall but retain at least one Type 3a interface with $\Delta G \sim -9$
that is nevertheless sufficient for recognition and for subsequent
filament disruption. All six protomer interfaces are required for
productive filament assembly; POPs need only one favourable recognition
surface to dock and at least one unfavourable surface to prevent the
downstream protomer from engaging, and that is enough to interrupt the
propagation of the assembly. Nonequilibrium assembly kinetics
additionally favour homotypic interactions unless POPs supply either
strongly unfavourable interfaces or multimeric steric interference, which
is consistent with the observation that POP2 and POP3 form elongated
oligomers after MBP cleavage \citep{mazanek2019fret}. The consistent
co-existence of rising FRET signal and absent mature filaments in nsEM
further suggests that POPs preferentially trap pseudo-nucleation
intermediates, stalling assembly before a stable filament is committed.

\subsection{Implications, limitations, and outlook}

The broader implication is that the same degenerate, code-like
recognition mechanism that allows chemically distinct receptor PYDs
(AIM2, IFI16, NLRP3, NLRP6) to converge on a single adaptor
ASC$_{PYD}$ also permits a small number of POPs to modulate a diverse
array of receptors \citep{lu2014unified, matyszewski2021distinct,
shen2019nlrp6, hochheiser2022directionality, sharif2019nek7,
andreeva2021nlrp3}. The picture is one of combinatorial interface
coding: target specificity is tuned not by one dominant sequence
contact but by the joint pattern of favourable and unfavourable
protomer-level interactions. This framework has obvious translational
appeal. Rationally engineered POP variants that accentuate
favourable-plus-unfavourable interface combinations at a specified
receptor could provide isoform-selective inflammasome inhibitors, and
small-molecule mimics that populate the same interface-energy
landscape could in principle achieve comparable selectivity.

Our study has several explicit limitations. NLRP2$_{PYD}$ did not form
filaments in HEK293T cells, which prevented a direct test of the
sequence-homology prediction for POP2 despite the 68\% identity.
IFI16$_{FL}$ is nuclear-localized \citep{kerur2011ifi16}, which
precludes direct co-imaging with cytosolic POPs in cells, and our
IC$_{50}$ rankings reflect in vitro solution conditions rather than the
full cellular context. The Rosetta interface analyses are likewise
static and do not capture conformational heterogeneity. Future work
should employ extended molecular dynamics and systematic mutagenesis
at favourable-plus-unfavourable interfaces to test the causal
contribution of each protomer contact, and should evaluate whether
engineered POP variants can selectively rescue or suppress
autoinflammatory phenotypes in transgenic models
\citep{sharma2018pyrin, vora2021inflammasome}.

In summary, this work reassigns POP1 from an ASC inhibitor to a decoy
adaptor for upstream receptors, identifies POP2 and POP3 as dual-target
nucleation-and-elongation blockers, and distills a protomer-interface
design principle for PYD filament inhibition. By replacing homology-based
guesswork with an interface-energy framework, our results open a
structurally grounded route to rational modulation of persistent
inflammasome activity in human disease.

\bibliography{references}

\end{document}
